\begin{abstract}
Este artículo presenta un estudio integral de factibilidad técnica y operativa para la implementación de un sistema de enlace satelital redundante basado en VENESAT-2. Se analizan los requisitos de disponibilidad de servicio superior al 99.9% en entornos tropicales, considerando arquitecturas de redundancia 1+1 y diversidad satelital/terrestre. Mediante simulaciones de balance de enlace en bandas Ku/Ka, evaluación de atenuación por lluvia utilizando modelos ITU-R y análisis FMEA de confiabilidad, se determinan parámetros óptimos de diseño. Los resultados indican viabilidad técnica con márgenes de enlace adecuados y mejoras operativas en resiliencia ante fallos. Se discuten implicaciones para la soberanía de comunicaciones en Venezuela y recomendaciones para implementación.
\end{abstract}